\documentclass[10pt,a4paper]{amsart}
\usepackage[margin=19mm]{geometry}
\usepackage{amsmath,amssymb,mathtools}
\usepackage[scr=rsfs,cal=euler]{mathalfa}
\usepackage{xfrac}
\usepackage[unicode,hidelinks,bookmarks=false]{hyperref}
\newcommand{\ie}{i.e.\ }
\newcommand{\cf}{cf.\ }
\newcommand{\resp}{respectively\ }
\newcommand{\Subsection}{\subsection}
\providecommand{\nobreakdash}{\nobreak\hbox{-}}
\setlength{\emergencystretch}{2em}
\multlinegap0pt
\usepackage{bm}
\usepackage{bbm}
\usepackage{dsfont}
\usepackage{xspace}
\usepackage{cancel}
\usepackage{mathtools}
\usepackage{amsthm}
\makeatletter
\@ifundefined{theorem}{\newtheorem{theorem}{Theorem}}{}
\@ifundefined{lemma}{\newtheorem{lemma}[theorem]{Lemma}}{}
\@ifundefined{proposition}{\newtheorem{proposition}[theorem]{Proposition}}{}
\makeatother
\newcommand{\ba}{\backslash}
\newcommand{\cC}{\mathcal{C}}
\title[Degeneracy: From Graphs to Matroids]{Degeneracy: From Graphs to Matroids}
\author{Allan Bickle, James Dylan Douthitt, Wayne Ge,, Jagdeep Singh}
\date{Source: arXiv:2608.09655v1 (CC BY 4.0)}
\begin{document}
\maketitle

\noindent\emph{Source and licence.} Degeneracy: From Graphs to Matroids. Version arXiv:2608.09655v1 (CC BY 4.0), distributed under the Creative Commons Attribution 4.0 International licence, as recorded in the arXiv licence metadata for this exact version and shown on the abstract page https://arxiv.org/abs/2608.09655v1. Full rights notice: licenses/arxiv-2608.09655-CC-BY-4.0.txt.\par\smallskip

\noindent\textit{Editorial scope.} Counted unit: the Lemma whose source label is \texttt{lem: 2-sum\_components\_extremal} in the article (source0.tex lines 343--346, source label \texttt{lem: 2-sum\_components\_extremal}), delivered together with the article's own complete original proof of it (source0.tex lines 348--419, one \texttt{proof} environment of 72 lines). No mathematics is written, repaired or re-derived: statement and proof are the article's own text line for line. Required context retained uncounted: prop: size\_bound (lines 155--158). Every definition, display and label of those lines is retained unchanged. Omitted from this package and not redistributed: 1--154, 159--342, 347--347, 420--761. Nothing inside a retained excerpt is omitted or altered: every retained body is delivered exactly as the article's lines are written, so no omission receipt of any kind accompanies this package. Citations: the retained excerpts carry no citation command at all, so no bibliography entry of the article is transcribed and no reference is redistributed. Reference adaptation: none was needed and none was made; every reference inside the delivered unit resolves inside it, so no unresolved reference or citation is produced. Printed slips kept literal: no printed word, symbol, hypothesis or number of the counted statement or of its proof was altered. Layout and class: the article's own class is \texttt{amsart} with packages including graphicx, amsthm, amsmath, amssymb, tikz, bm, mathtools, xifthen, comment, geometry, enumitem, todonotes, hyperref; the wrapper is the lane's pinned \texttt{amsart} editorial wrapper at 10pt on A4, which restates the theorem-like environments the retained text needs and adds the article's own macro definitions that the retained text needs (\texttt{\textbackslash ba}, \texttt{\textbackslash cC}), with extra wrapper packages (bm, bbm, dsfont, xspace, cancel, mathtools). The delivered numbering is the unit's own. Assets: no retained span contains a figure environment, a graphics-inclusion command, a PDF-page inclusion, a table or a TikZ picture; no image file is copied into this package, the package declares no asset dependency and no image file is redistributed with it. Licence: version arXiv:2608.09655v1 of this article is distributed under the Creative Commons Attribution 4.0 International licence, as recorded in the arXiv licence metadata for this exact version and shown on the abstract page https://arxiv.org/abs/2608.09655v1. A full rights notice is delivered at licenses/arxiv-2608.09655-CC-BY-4.0.txt. Characters: every retained excerpt is pure ASCII, so the delivered engine, XeLaTeX with its default Latin Modern text font, reports no missing character.
\begin{proposition}
\label{prop: size_bound}
Let $M$ be a $k$-degenerate matroid of rank $r$. Then $|E(M)|\leq kr$. If $M$ is simple, then $|E(M)| \leq k(r-1)+1$. 
\end{proposition}

\begin{lemma}
\label{lem: 2-sum_components_extremal}
For $k \ge 3$, suppose $M = M_1 \oplus_2 M_2$ at basepoint $p$, where $M$ is a simple minimally $k$-degenerate matroid of extremal size $|E(M)| = (k-1)(r(M)-1)+2$. Then both $M_1$ and $M_2$ are simple minimally $k$-degenerate matroids of extremal size.
\end{lemma}

\begin{proof}
Let
\[
    f(s)=(k-1)(s-1)+2.
\]
For each $i\in\{1,2\}$, the matroid $M_i\ba p$ is a restriction of $M$, and hence is simple and $(k-1)$-degenerate. Moreover, at least one of $M_1$ and $M_2$ is simple. Indeed, since $M_i\ba p$ is simple, every parallel pair of $M_i$ must contain $p$. If both $M_1$ and $M_2$ had an element parallel to $p$, then those two elements would be parallel in $M$, contradicting the simplicity of $M$. Without loss of generality, we may assume that $M_1$ is simple.

Suppose that $M_1$ is $(k-1)$-degenerate. By Proposition~\ref{prop: size_bound},
\[
    |E(M_1)|\leq (k-1)(r(M_1)-1)+1=f(r(M_1))-1.
\]
Since $p$ is not a coloop of $M_2$, we have $r(M_2\ba p)=r(M_2)$. Applying Proposition~\ref{prop: size_bound} to the simple $(k-1)$-degenerate matroid $M_2\ba p$ gives
\[
    |E(M_2)|-1\leq (k-1)(r(M_2)-1)+1,
\]
and hence $|E(M_2)|\leq f(r(M_2))$. Therefore,
\begin{align*}
    |E(M)|
    &=|E(M_1)|+|E(M_2)|-2\\
    &\leq f(r(M_1))+f(r(M_2))-3\\
    &=f(r(M))-1,
\end{align*}
contradicting $|E(M)|=f(r(M))$. Thus, $M_1$ is not $(k-1)$-degenerate. Consequently, $M_1$ has a restriction $N_1$ such that $g^*(N_1)\geq k$. Since $M_1\ba p$ is $(k-1)$-degenerate, we have $p\in E(N_1)$.

We next show that $M_2$ is simple. Suppose otherwise. Then there is an element $e\in E(M_2)-\{p\}$ that is parallel to $p$ in $M_2$. The matroid $M$ has a proper restriction isomorphic to the matroid obtained from $N_1$ by relabeling $p$ as $e$. This restriction has cogirth at least $k$, contradicting the minimal $k$-degeneracy of $M$. Therefore, $M_2$ is simple.

By symmetry, the preceding counting argument shows that $M_2$ is not $(k-1)$-degenerate. Hence $M_2$ has a restriction $N_2$ such that $g^*(N_2)\geq k$, and necessarily $p\in E(N_2)$. Since $g^*(N_i) \geq k$ for each $i \in \{1,2\}$, the basepoint $p$ is not a coloop of $N_i$. Therefore, $N_1\oplus_2 N_2$ is well-defined.

Recall that the cocircuits of $N_1\oplus_2 N_2$ are the members of
\begin{align*}
    &\cC^*(N_1/p)\cup\cC^*(N_2/p)\\
    &\quad\cup\bigl\{(C_1^*\cup C_2^*)\backslash\{p\}:
    p\in C_1^*\in\cC^*(N_1)\ \text{and}\
    p\in C_2^*\in\cC^*(N_2)\bigr\}.
\end{align*}
Every member of $\cC^*(N_1/p)\cup\cC^*(N_2/p)$ has size at least $k$. Moreover, if $p\in C_i^*\in\cC^*(N_i)$ for each $i\in\{1,2\}$, then
\[
    \bigl|(C_1^*\cup C_2^*)\backslash\{p\}\bigr|
    =|C_1^*\backslash\{p\}|+|C_2^*\backslash\{p\}|
    \geq 2k-2\geq k.
\]
Therefore, $g^*(N_1\oplus_2 N_2)\geq k$. Since $N_1\oplus_2 N_2$ is a restriction of $M$, the minimality of $M$ implies that $N_1=M_1$ and $N_2=M_2$. In particular, $M_1$ and $M_2$ are simple matroids with cogirth at least $k$.

We now show that $M_1$ is minimally $k$-degenerate. Suppose that $R_1$ is a nonempty proper restriction of $M_1$ such that $g^*(R_1)\geq k$. If $p\notin E(R_1)$, then $R_1$ is a proper restriction of $M$, a contradiction. Thus, $p\in E(R_1)$. Since $g^*(R_1)\geq k$, the element $p$ is not a coloop of $R_1$, so $R_1\oplus_2 M_2$ is well-defined. The cocircuits of $R_1\oplus_2 M_2$ are the members of
\begin{align*}
    &\cC^*(R_1/p)\cup\cC^*(M_2/p)\\
    &\quad\cup\bigl\{(C_1^*\cup C_2^*)\backslash\{p\}:
    p\in C_1^*\in\cC^*(R_1)\ \text{and}\
    p\in C_2^*\in\cC^*(M_2)\bigr\}.
\end{align*}
Every member of the first two families has size at least $k$, and every member of the third family has size at least $2k-2$. Therefore, $g^*(R_1\oplus_2 M_2)\geq k$. However, $R_1\oplus_2 M_2$ is a proper restriction of $M$, a contradiction. Therefore, $M_1$ is minimally $k$-degenerate. By symmetry, $M_2$ is also minimally $k$-degenerate.

Applying the size bound to $M_1$ and $M_2$ gives $|E(M_i)|\leq f(r(M_i))$ for each $i\in\{1,2\}$. Since
\[
    r(M)=r(M_1)+r(M_2)-1
\]
and
\[
    f(r(M))=f(r(M_1))+f(r(M_2))-2,
\]
we have
\begin{align*}
    f(r(M_1))+f(r(M_2))-2
    &=|E(M)|\\
    &=|E(M_1)|+|E(M_2)|-2.
\end{align*}
It follows that
\[
    |E(M_i)|=f(r(M_i))=(k-1)(r(M_i)-1)+2
\]
for each $i\in\{1,2\}$. Thus, both $M_1$ and $M_2$ are simple minimally $k$-degenerate matroids of extremal size.
\end{proof}

\end{document}
